% Generated by Listas/make_statements.py from pratica-lecture-04.md, item 5. Do not edit.
\textbf{Source.} Tadelis (2013), \emph{Game Theory: An Introduction}, Exercise 8.13 ``Junk Mail Advertising'', pp.~173--174 {[}PDF 192--193{]}. Reproduced verbatim. Answer in the source: none in the book (the preface, p.~xiii, says about half are solved on Princeton UP's website; not checked).



\begin{quote}

\textbf{Junk Mail Advertising:} Suppose a single good is owned by a single seller who values it at \(c > 0\) (he can consume the good and get a payoff of \(c\)). There is a single buyer who has a small transportation cost \(k > 0\) to get to and back from the seller's store, and he values the good at \(v > c + k\). The buyer first decides whether to make the commute or stay at home, not buy the good, and receive a payoff of 0. If the buyer commutes to the store, the seller can then make the buyer a take-it-or-leave-it price offer \(p \geq 0\). The buyer can then accept the offer, pay \(p\), and get the good, or he can walk out and not buy the good. Assume that \(c\), \(v\), and \(k\) are common knowledge.



\begin{enumerate}

\def\labelenumi{\alph{enumi}.}

\item

  As best you can, draw the extensive form of this game. What is the best response of the buyer at the node at which he decides whether to accept or reject the seller's offer?

\item

  Find the subgame-perfect equilibrium of the game and show that it is unique. Is it Pareto optimal?

\item

  Are there Nash equilibria that yield a higher payoff to both players as compared to the subgame-perfect equilibrium you found in (b)?

\item

  Now assume that before the game is played the seller can, at a small cost \(\varepsilon < (v - c - k)\), send the buyer a postcard that commits the seller to a certain price at which the buyer can buy the good (e.g., ``bring this coupon and get the good at a price \(p\)''). Would the seller choose to do so? Justify your answer with an equilibrium analysis.

\end{enumerate}

\end{quote}

